% Generated draft: requires mathematical and source review.
\documentclass[11pt]{article}
\usepackage[margin=25mm]{geometry}
\usepackage{fontspec}
\setmainfont{Latin Modern Roman}
\usepackage{amsmath,amssymb,mathtools,unicode-math}
\setmathfont{Latin Modern Math}
\usepackage{enumitem,needspace,xurl,hyperref}
\setlength{\parindent}{0pt}
\setlength{\parskip}{5pt}
\setlength{\emergencystretch}{3em}
\newcommand{\N}{\mathbb{N}}
\newcommand{\Z}{\mathbb{Z}}
\newcommand{\Q}{\mathbb{Q}}
\newcommand{\R}{\mathbb{R}}
\newcommand{\C}{\mathbb{C}}
\newcommand{\F}{\mathbb{F}}
\newcommand{\A}{\mathbb{A}}
\newcommand{\PP}{\mathbb{P}}
\newcommand{\E}{\mathbb{E}}
\newcommand{\eps}{\varepsilon}
\newcommand{\dd}{\,\mathrm{d}}
\begin{document}
\section*{UnsolvedMath ID: 30001345}
\textbf{Source problem code:} catalogue entry
\subsection{Definitions and mathematical statement}
The supplied record does not define the \(M\)-number or specify the category of deformations. Consequently, the following is a conditional formulation of the question; it does not supply the missing definition of \(M\).

Let \((T,t_0)\) be a pointed parameter space for an admissible deformation \((C_t)_{t\in T}\) of plane-curve singularities, with special member \(C_{t_0}\). Here the topology on \(T\), admissibility of the deformation, and the meaning of its singularity members must be specified from the source. Write \(M(C_t)\) for the intended numerical invariant. If this invariant is real-valued, upper semicontinuity at \(t_0\) means
\[
\forall\varepsilon>0\quad\exists U\text{ an open neighborhood of }t_0\quad\forall t\in U,\qquad M(C_t)<M(C_{t_0})+\varepsilon.
\]
The general question asks whether this condition holds at every parameter value in every admissible deformation. The particular question asks whether it holds whenever the deformation is \(\delta\)-constant, meaning that, on a neighborhood \(V\) of \(t_0\),
\[
\delta(C_t)=\delta(C_{t_0})\qquad\text{for every }t\in V.
\]

Under the conventional interpretation in terms of reduced complex analytic plane-curve germs \((C_t,p_t)\), the local \(\delta\)-invariant is
\[
\delta(C_t,p_t)=\dim_{\C}\!\left(\widetilde{\mathcal O}_{C_t,p_t}/\mathcal O_{C_t,p_t}\right),
\]
where \(\widetilde{\mathcal O}_{C_t,p_t}\) is the integral closure of the local ring in its total ring of fractions. This interpretation is conditional: the supplied record does not state reducedness, the ground field, or an analytic category. It also does not say whether invariants are taken at a selected singular point or summed over singular points appearing in a nearby fiber. Thus the displayed formulation requires these choices, together with the definition and normalization of \(M\), before it becomes a fully specified mathematical problem.

The supplied dated literature assessment says that upper semicontinuity fails without the \(\delta\)-constant hypothesis, while only partial bounds are known for the \(\delta\)-constant question and that question remains unresolved. These are supplied curation claims, not independently verified conclusions of this draft.
\subsection{Short English statement}
The question asks whether the \(M\)-number can increase near a special plane-curve singularity under deformation, especially when the \(\delta\)-invariant stays constant. The supplied assessment reports failure in general and an unresolved \(\delta\)-constant case with partial bounds. A precise unconditional statement still requires a definition of \(M\), deformation hypotheses, and a convention for fibers with several singular points.
\subsection{Sources}
\begin{itemize}
\item Supplied local UnsolvedMath registry record, permanent ID \(30001345\), source problem code \texttt{catalogue entry}, in the collection \emph{Oberwolfach Reports Open Problems}. No individual proposer is identified; the supplied proposed year is 2009.
\item Supplied source citation: \emph{Singularities (2010)}, Oberwolfach Reports, pp.\ 2405--2470, DOI \texttt{10.4171/owr/2009/43}. Original report: \url{https://doi.org/10.4171/owr/2009/43}. The record describes these reports as documenting mathematical workshops and their open-problem discussions.
\item Supplied literature reference, labeled ``deformations of plane curve singularities'': \url{https://arxiv.org/abs/0907.4129}.
\item Supplied literature reference, labeled ``official publication record'': \url{https://doi.org/10.18910/50809}.
\item Literature assessment dated exactly \texttt{2026-08-21}: the supplied curation labels the problem partially solved, reports failure of upper semicontinuity without the \(\delta\)-constant hypothesis, and describes the \(\delta\)-constant case as unresolved with partial bounds. This is supplied dated curation derived from the OWR material, not an independently verified literature assessment; the cited claims and time-sensitive status require independent verification.
\item The supplied statement label is \texttt{corrected\_verified}; its accompanying verification describes meaning-preserving semantic, typographic, and LaTeX cleanup without added source-dependent content. This editorial label does not establish mathematical correctness or resolve the missing definitions.
\end{itemize}
\end{document}
